\section{Exact executable certificate}\label{sec:certificate}

The accompanying executable artifact checks the convention-sensitive finite
algebra using Gaussian numbers represented as pairs of
\texttt{fractions.Fraction} values.  It evaluates the four energies, performs
the $C_4$ transform, recomputes both sides of the defect identity, and checks
exact witnesses at the center, interior, and boundary of the feasible disk.
No floating-point comparison appears in a theorem-facing check.

The producer has a required mutually exclusive interface: \texttt{--write}
creates the canonical certificate, while \texttt{--check} is read-only.  The
JSON representation uses sorted keys, compact separators, ASCII encoding,
one terminal newline, reduced fractions with positive denominators, and a
SHA-256 digest over the payload.  A separately implemented checker parses the
document with duplicate-key and nonfinite-constant rejection and recomputes
the exact fixtures without importing the producer.

Both programs exercise semantic mutation firewalls.  They reject an
inner-product orientation reversal, replacement of $\ii^j$ by $(-\ii)^j$,
a nonzero $F_2$, leakage of a common offset into $F_1$, a fabricated physical
attachment, a status promotion, and Python's boolean-as-integer ambiguity.
The digest is recomputed on semantic mutations before validation, so mutation
rejection cannot be credited merely to stale checksums.  Parser mutations
also test duplicate keys and nonfinite constants.

Finally, an independent stress program exhausts all ordered pairs of
two-dimensional vectors whose component real and imaginary parts lie in
$\{-1,0,1\}$.  This is $81^2=6561$ vector pairs and $26244$ phase-energy
evaluations.  Every calculation is exact, but the census is labelled
\begin{center}
\texttt{NUMERICAL\_FINITE\_ILLUSTRATION\_ONLY}.
\end{center}
The executable artifacts test transcription, representation, and finite
interfaces.  The symbolic proofs in Sections~\ref{sec:spectrum}
and~\ref{sec:disk}, rather than the finite census, establish the theorem.
